\section{Introdução}\label{sec:intro}

O federalismo fiscal brasileiro atribui aos municípios uma parcela expressiva da provisão de serviços públicos (educação básica, atenção primária à saúde, mobilidade e zeladoria urbana), mas reserva a eles uma base tributária estreita, concentrada no ISS, no IPTU e no ITBI\footnote{BRASIL. \textit{Constituição da República Federativa do Brasil de 1988}. Disponível em: \url{https://www.planalto.gov.br/ccivil_03/constituicao/constituicao.htm}.}. O hiato entre responsabilidades de gasto e capacidade de arrecadação --- o desequilíbrio fiscal vertical --- é coberto por transferências constitucionais e legais da União e dos estados, como o Fundo de Participação dos Municípios (FPM), a cota-parte do ICMS, o FUNDEB e os repasses fundo a fundo do SUS. A teoria do federalismo fiscal reconhece o papel dessas transferências na correção de externalidades e na equalização de capacidades \cite{oates1999,bird2002}, mas a chamada segunda geração da literatura chama atenção para os incentivos que elas criam: governos que gastam recursos arrecadados por outros tendem a perceber o custo do gasto de modo distorcido, a arrecadar menos e a pressionar por novos repasses \cite{weingast2009,rodden2002,oates2008}.

Dois fenômenos concentram esse debate. O primeiro é o efeito \textit{flypaper}: uma transferência incondicional eleva o gasto local muito mais do que um aumento equivalente da renda dos residentes, ao contrário do que prevê o teorema de equivalência derivado do modelo do eleitor mediano \cite{bradford1971,hines1995}. O segundo é a ilusão fiscal, isto é, a percepção sistematicamente distorcida do custo dos bens públicos pelos contribuintes, que pode ser reforçada quando o financiamento vem de outra esfera de governo \cite{wagner1976,logan1986,dollery1996}. No Brasil, a literatura documenta ambos os fenômenos para o conjunto dos municípios \cite{cossio2001,mattos2011,ferreira2019,ferreira2022}, e a tese de Ferreira (2021)\footnote{FERREIRA, L. O. G. \textit{Constitutional tax transfer in Brazilian municipalities: the fiscal illusion and the flypaper effect}. 2021. Tese (Doutorado) -- Universidade de Brasília, Brasília, 2021.}, referência de partida deste trabalho, associa a magnitude do efeito \textit{flypaper} à ilusão fiscal gerada pelas transferências constitucionais.

A Amazônia oferece um contraste particularmente informativo. Manaus, capital do Amazonas e sede da Zona Franca, concentra a atividade industrial do estado e, por isso, tem base própria de ISS e uma cota-parte de ICMS relativamente elevada \cite{serafico2005}. Parintins, município de porte médio situado em uma ilha do baixo Amazonas e conhecido nacionalmente pelo Festival Folclórico, tem PIB per capita cerca de um quarto do de Manaus e depende quase integralmente de recursos de terceiros. As duas cidades compartilham o mesmo estado, as mesmas regras de repartição e o mesmo tribunal de contas, mas diferem radicalmente em escala e em estrutura produtiva, o que permite observar como a dependência de transferências se manifesta em contextos econômicos distintos.

Este trabalho atualiza e amplia a base construída na prova anterior atribuída a esta matrícula\footnote{VICTOR, C. \textit{Prova P1 de Finanças do Setor Público} (matrícula 231020069). 2026. Repositório GitHub: \url{https://github.com/cauaviictor/P1---FSP---231020069---CAUA-VICTOR}, \textit{commit} 5209bcec.}, que analisava a Espanha, Maceió/AL e Duque de Caxias/RJ no período 2013--2022. Com a redistribuição dos conjuntos feita pelo professor, os municípios passaram a ser Manaus e Parintins, mantendo-se a Espanha como país de comparação. A pergunta de pesquisa é: \textbf{em que medida as transferências intergovernamentais condicionam o financiamento e o gasto de Manaus e Parintins entre 2016 e 2025, e esse padrão é compatível com o efeito \textit{flypaper} e com a hipótese de ilusão fiscal, quando confrontado com a organização fiscal subnacional da Espanha?}

O objetivo geral é caracterizar a dependência de transferências dos dois municípios e estimar, com os dados anuais desses dois municípios, a associação entre transferências, renda local e gasto. Os objetivos específicos são: (i) reconstruir, documentar e corrigir a base de dados do trabalho anterior a partir de fontes oficiais, com valores reais e per capita; (ii) mapear a literatura sobre transferências, efeito \textit{flypaper} e ilusão fiscal por meio de um estudo bibliométrico; (iii) descrever a evolução das receitas, das despesas e dos indicadores de autonomia de Manaus e Parintins; (iv) testar se o gasto municipal responde mais às transferências do que à renda local e se as transferências se associam a menor esforço tributário; e (v) contrastar esses padrões com os das Comunidades Autônomas e das entidades locais espanholas.

A justificativa é dupla. Do ponto de vista substantivo, a Reforma Tributária aprovada pela Emenda Constitucional n.º 132/2023 vai substituir o ISS e o ICMS pelo IBS ao longo da próxima década\footnote{BRASIL. Emenda Constitucional n.º 132, de 20 de dezembro de 2023 (Reforma Tributária). Disponível em: \url{https://www.planalto.gov.br/ccivil_03/constituicao/emendas/emc/emc132.htm}.}, alterando justamente as bases que distinguem Manaus de Parintins; uma linha de base bem documentada da dependência atual é condição para avaliar essa transição. Do ponto de vista metodológico, os trabalhos anteriores da disciplina utilizaram planilhas extraídas manualmente, sem registro das classificações contábeis, dos deflatores e das quebras de série; a reconstrução reprodutível aqui proposta permite conferir cada número a partir das fontes. A contribuição do estudo está em analisar em detalhe dois municípios amazônicos raramente estudados, tratar explicitamente as mudanças do plano de contas e da base populacional e situar os resultados em perspectiva internacional com dados do Sistema Europeu de Contas.

A relação entre a análise municipal e a análise do país não pretende comparar Manaus com a Espanha como unidades equivalentes. A Espanha é um Estado unitário fortemente descentralizado, organizado em Comunidades Autônomas e entidades locais\footnote{ESPAÑA. \textit{Constitución Española de 1978}. Disponível em: \url{https://www.boe.es/eli/es/c/1978/12/27/(1)/con}.}, cujo financiamento subnacional combina tributos próprios, impostos cedidos e transferências de nivelamento \cite{zabalza2011,toboso2010}. O país serve, assim, como referência institucional: permite observar como níveis subnacionais com maior autonomia tributária e sujeitos a regras fiscais nacionais e europeias se comportam em termos de dependência, resultado e endividamento. A comparação é feita em proporções da receita e do PIB, que neutralizam as diferenças de moeda e escala, e por subsetor das administrações públicas, aproximando o nível de governo analisado.

O artigo está organizado em quatro seções além desta introdução. A Seção~\ref{sec:referencial} apresenta o estudo bibliométrico e discute a literatura. A Seção~\ref{sec:metodo} descreve o delineamento, as fontes, as variáveis, os tratamentos e as técnicas de análise. A Seção~\ref{sec:resultados} apresenta as mudanças em relação ao trabalho anterior, as estatísticas descritivas, as estimações e as limitações. A Seção~\ref{sec:conclusao} conclui.
